\documentclass{article}

\textwidth=6.5in
\textheight=8.5in
\voffset=-54pt
\oddsidemargin=0in
\evensidemargin=0in
\setlength{\parskip}{8pt}
\setlength{\parindent}{0pt}

\usepackage{workbook}

\usepackage[sols]{handout}

\begin{document}

\begin{center}
  \large\textsc{Problem Set 29 - Integration by Parts}
  
      \pspace{12pt}
  \begin{problemonly}\large Name: \rule{5cm}{0.15mm}\\ \end{problemonly}
\end{center}

\begin{grading}
    \begin{enumerate}
        \item 11 - 2/2/4/3
        \item 12 - 3/4/5
        \item 10 - 5/5
        \item 4
    \end{enumerate}
Total: 37 points
\end{grading}

\begin{enumerate}

\begin{grading}
Parts (a) and (bi) are direct applications of integration by parts. Part (bii) and (c) involve two applications but utilize previous computations.
\end{grading}


\item Find each indefinite integral using integration by parts.

\begin{enumerate}

\item 
\begin{problem}[\vfill]
    $\displaystyle \int x \cos(x)  \,dx $
\end{problem}

\begin{solution}
Set $u = x$ and $\frac{dv}{dx} = \cos(x)$. Then, $\frac{du}{dx} = 1$ and $v = \sin(x)$. Therefore,
\[\int x \cos(x) \,dx = x\sin(x) - \int \sin(x) \,dx = x\sin(x) + \cos(x) + C.\]
\end{solution}

    \item 
    \begin{problem}[\vfill] 
    $\displaystyle \int \ln(x)\, dx$
    \end{problem}

    \begin{solution}
    
    Let $u = \ln x$ and $\frac{dv}{dx} = 1$. 
    Then $\frac{du}{dx} = \frac{1}{x}$ and $v = x$.
    Using the integration by parts formula,
    \begin{align*}
    \int \ln(x) \, dx &= x\ln(x) - \int \frac{x}{x}\, dx \\
    &= x\ln(x) - x + C.
    \end{align*}
        
    \end{solution}

    \item
    \begin{problem}[\vfill\newpage]
    $\displaystyle \int (\ln x)^2 \, dx$
    \end{problem}

    \begin{solution}

    Let $u = \ln x$ and $\frac{dv}{dx} = \ln x$. Then
    $\frac{du}{dx} = \frac{1}{x}$ and $v = x\ln(x) - x$.
    Using the integration by parts formula,

    \begin{align*}
        \int (\ln x)^2 \, dx &=
        \ln(x)\left[x\ln(x)-x\right] - \int \frac{x\ln(x)-x}{x} \, dx \\
        &= x(\ln(x))^2 - x\ln(x) - \int (\ln(x) - 1)\, dx \\
        &= x(\ln(x))^2 - x\ln(x) - \int \ln(x)\, dx + \int 1\,dx\\
        &= x(\ln(x))^2 - x\ln(x) - [x\ln(x)-x] + x + C \\
        &= x(\ln(x))^2 - 2x\ln(x) + 2x + C
    \end{align*}

    \end{solution}


\item  
\begin{problem}[\vfill]
    $\displaystyle \int \arctan(x) \,dx $
\end{problem}

\begin{solution}
Let $u=\arctan(x)$ and $dv = dx$. Then, $du = \dfrac{1}{1+x^2}dx$ and $v=x$, 
\begin{align*}
\int \arctan(x)\,dx &= x\arctan(x) - \int \dfrac{x}{1+x^2}\,dx \\
\intertext{The second integral can be done using substitution,}
&= x\arctan(x) - \frac{1}{2}\ln|1+x^2| +C
\end{align*}
\end{solution}

\item  
\begin{problem}[\vfill]
    $\displaystyle \int x^2 \sin(x) \,dx $
\end{problem}

\begin{solution}
Using integration by parts twice, first let $u = x^2$ and $dv = \sin(x)dx$. Then, 
\[\int x^2 \sin(x) \,dx = -x^2\cos(x) + 2 \int x\cos(x) \,dx.\]
The remaining integral $\int x\cos(x) \,dx$ is solved as in part (a). So, 

$$\displaystyle \int x^2 \sin(x) \,dx = -x^2\cos(x) + 2 ( x\sin(x) + \cos(x))$$
\end{solution}

\item  
\begin{problem}[\vfill]
    $\displaystyle \int x^3 e^{2x} \,dx $
\end{problem}

\begin{solution}
We need to use integration by parts three times. Let $u=x^3$ and $dv=e^{2x}dx$. Then $du = 3x^2dx$ and $v=\tfrac{1}{2}e^{2x}$.
\begin{align*}
    \int x^3 e^{2x}\,dx &= \tfrac{1}{2}x^3e^{2x} - \tfrac{3}{2}\!\int x^2e^{2x}\,dx 
    \intertext{Now we use integration by parts again on the second integral with $u=x^2$ and $dv = e^{2x}dx$. Then $du = 2xdx$ and $v=\tfrac{1}{2}e^{2x}$.}
    &= \tfrac{1}{2}x^3e^{2x} - \tfrac{3}{2}\!\left(\tfrac{1}{2}x^2e^{2x}-\!\int xe^{2x}\,dx\right) \\
    &= \tfrac{1}{2}x^3e^{2x} -\tfrac{3}{4}x^2e^{2x} + \tfrac{3}{2}\!\int xe^{2x}\,dx
    \intertext{Now we do integration by parts one last time on the last integral with $u=x$ and $dv=e^{2x}dx$, so $du=dx$ and $v=\tfrac{1}{2}e^{2x}$.}
    &= \tfrac{1}{2}x^3e^{2x} -\tfrac{3}{4}x^2e^{2x} + \tfrac{3}{2}\!\left(\tfrac{1}{2}xe^{2x}-\tfrac{1}{2}\!\int e^{2x}\,dx\right)\\
    &= \tfrac{1}{2}x^3e^{2x} -\tfrac{3}{4}x^2e^{2x} + \tfrac{3}{4}xe^{2x}-\tfrac{3}{4}\!\int e^{2x}\,dx\\
    \intertext{This last integral we do not need integration by parts to do.}
        &= \tfrac{1}{2}x^3e^{2x} -\tfrac{3}{4}x^2e^{2x} + \tfrac{3}{4}xe^{2x}-\tfrac{3}{8}e^{2x}+C \\
        &= \tfrac{1}{8}e^{2x} (4x^3-6x^2+6x-3) + C\\
\end{align*}
If you need to integrate by parts repeatedly, there is a method known as the ``tabular method" or the ``DI method" that makes the process more efficient. You are encouraged to ask a TF or research it if you are interested, but we will not require you to know it.
\end{solution}

\end{enumerate}

\pspace{\newpage}

\item 
\begin{grading}
    Part (a) is a direct application of integration by parts with an evaluation of the definite integral (2 points for the IBP and 1 point for the evaluation). Part (b) involves first integration by parts and then integration by substitution (with 2 points for the IBP and 2 points for the substitution) 
\end{grading}

Find or evaluate the following integrals. You should show your work and carefully justify your answers. In particular, if you use $u$-substitution, be sure to state what $u$ and $du$ are; if you use integration by parts, be sure to state what $u$, $dv$, $du$, and $v$ are.

\begin{enumerate}

\item
\begin{problem}[\vfill]
    $\displaystyle \int_1^e x^2 \ln(x) \, dx$
\end{problem}

\begin{solution}
    This integral can be solved by integration by parts where
    \begin{align*} 
    u = \ln x \qquad &dv = x^2 \, dx \\
    du = \frac{1}{x} \, dx \qquad &v = \frac{1}{3}x^3
    \end{align*}

Solving first for the indefinite integral, we have
\begin{align*}
    \int x^2 \ln(x)\, dx &= 
    \frac{1}{3}x^3 \ln(x) - \int \frac{1}{3}x^3 \cdot \frac{1}{x}\, dx \\
    &= \frac{1}{3}x^3 \ln(x) - \frac{1}{3} \int x^2\, dx \\
    &= \frac{1}{3}x^3 \ln(x) - \frac{1}{9}x^3 + C.
\end{align*}

Now we are ready to solve the definite integral:
\[\int_1^e x^2 \ln(x)\, dx = 
\left[\tfrac{1}{3}x^3 \ln(x) - \tfrac{1}{9}x^3\right]_1^e
= \tfrac{1}{3}e^3 \cdot 1 - \tfrac{1}{9}e^3 -(\tfrac{1}{3}\cdot 0-\tfrac{1}{9}) = \tfrac{1}{3}e^3 - \tfrac{1}{9} (e^3 - 1).\]
\end{solution}

\item
\begin{problem}[\vfill]
    $\displaystyle \int x^3 \cos(x^2) \, dx$
\end{problem}

\begin{solution}
    Integration by parts works here. You may be tempted to
    use $u = x^3$ and $dv = \cos(x^2)\, dx$, however it's not
    clear how to find the antiderivative of $\cos(x^2)$. Instead
    we use integration by parts where,
    \begin{align*}
        u = x^2 \qquad & dv = x\cos(x^2) \, dx \\
        du = 2x \, dx \qquad & v = \frac{1}{2}\sin(x^2).
    \end{align*}

    We have that
    \[\int x^3\cos(x^2) = \frac{1}{2}x^2\sin(x^2) -
    \int x\sin(x^2) \, dx,\]

    we can solve the remaining integral with the substitution
    $w = x^2$, $dw = 2x\,dx$:
    \begin{align*}
    \int x^3 \cos(x^2)\, dx &= \frac{1}{2}x^2 \sin(x^2) - \frac{1}{2} \int \sin(w)\,dw \\
    &= \frac{1}{2}x^2 \sin(x^2) + \frac{1}{2} \cos(x^2) + C.
    \end{align*}
\end{solution}

\item
\begin{problem}[\vfill\newpage]
    $\displaystyle \int_0^1 x^2 e^{3x} \,dx$
\end{problem}

\begin{solution}
    We start with integration by parts with
    \begin{align*}
        u = x^2 \qquad & dv = e^{3x} \, dx \\
        du = 2x \, dx \qquad & v = \frac{1}{3}e^{3x}.
    \end{align*}
    Using the integration by parts formula:
    \[\int x^2 e^{3x}\, dx = \frac{1}{3}x^2e^{3x} 
    - \frac{2}{3} \int xe^{3x}\, dx. \]

    We need to use integration by parts again, with
    \begin{align*}
        u = x \qquad & dv = e^{3x} \, dx \\
        du = 1 \, dx \qquad & v = \frac{1}{3}e^{3x}.
    \end{align*}
    Substituting the integration by parts formula:
    \begin{align*}
    \int x^2 e^{3x}\, dx &= \frac{1}{3}x^2e^{3x} 
    - \frac{2}{3} \int xe^{3x}\, dx \\
    &= \frac{1}{3}x^2e^{3x} 
    - \frac{2}{3} \left[\frac{1}{3}xe^{3x} - 
    \frac{1}{3}\int e^{3x}\, dx \right] \\
    &= \frac{1}{3}x^2e^{3x} 
    - \frac{2}{3} \left[\frac{1}{3}xe^{3x} - 
    \frac{1}{9} e^{3x} \right] \\
    &= \frac{1}{3}x^2e^{3x} - \frac{2}{9}x e^{3x} + \frac{2}{27}e^{3x} + C.
    \end{align*}
    Now we can evaluate the definite integral:
    \[\int_0^1 x^2 e^{3x}\, dx 
    = \left[\tfrac{1}{3}x^2e^{3x} - \tfrac{2}{9}x e^{3x} + \tfrac{2}{27}e^{3x}\right]_0^1 
    = \tfrac{1}{3}e^3 - \tfrac{2}{9}e^3 + \tfrac{2}{27}e^3 - (0 - 0 + \tfrac{2}{27}) = \tfrac{5}{27}e^3 -\tfrac{2}{27}
    \]
\end{solution}

\end{enumerate}

\item 
\begin{grading}
    These both involve substitution and integration by parts. The substitution needs to happen first to get to IBP.
\end{grading}

Find each indefinite integral. \textit{Hint: Try a substitution first!}

\begin{enumerate}

\item
\begin{problem}[\vfill]
    $\displaystyle \int x \ln(x^2+1) \, dx$
\end{problem}

\begin{solution}
    We can use the substitution $w=x^2+1$, $dw = 2x\, dx$.
    \begin{align*}
        \int x \ln(x^2+1)\, dx &= \frac{1}{2}\int \ln w \, dw  
    \end{align*}
    Now we can use integration by parts with
    \begin{align*}
        u = \ln w \qquad & dv = dw \\
        du = \frac{1}{w}\, dw \qquad & v = w
    \end{align*}
    giving us
    \begin{align*}
        \int x \ln(x^2+1)\, dx &= \frac{1}{2}\int \ln w \, dw  \\
        &= \frac{1}{2}\left[w\ln w - \int 1 \, dw \right] \\
        &= \frac{1}{2}w\ln(w) - \frac{1}{2}w + C\\
        &= \frac{1}{2}(x^2+1)\ln(x^2+1) -\frac{1}{2}(x^2+1) +C. \\
    \end{align*}
\end{solution}

\item 
\begin{problem}[\vfill\newpage]
    $\displaystyle \int -\frac{\ln(x+1)}{(x+1)^2} \, dx$
\end{problem}

\begin{solution}
    We can use the substitution $w = x+1$, giving us
    \[\int -\frac{\ln(x+1)}{(x+1)^2} \, dx = -\int \frac{\ln(w)}{w^2} \, dw \]
    Then we can do integration by parts with
    \begin{align*}
        u = \ln(w) \qquad & dv = \frac{1}{w^2}\, dw \\
        du = \frac{1}{w}\,dw \qquad & v = -\frac{1}{w}
    \end{align*}
    giving us
    \begin{align*}
    \int -\frac{\ln(x+1)}{(x+1)^2} \, dx &= -\int \frac{\ln(w)}{w^2} \, dw\\
     &=  -\left[-\frac{\ln(w)}{w} + \int \frac{1}{w^2}\,dw \right] \\
     &= \frac{\ln(w)}{w} + \frac{1}{w} + C \\
     &= \frac{\ln(x+1)}{x+1} + \frac{1}{x+1} + C \\
     &= \frac{\ln(x+1)+1}{x+1} + C.
    \end{align*}
\end{solution}

\end{enumerate}



\item

\begin{problem}[\vfill]
    Consider $\displaystyle \int x^{100}e^x\,dx$. Which function, $x^{100}$ or $e^x$, is a better choice for $u$ in the integration by parts formula? How many times would you need to apply the formula to get a final answer? Explain.
\end{problem}

\begin{solution}
    It is a better choice to let $u=x^{100}$. If we let $u=e^x$, then we get a more complicated integral than the one we started with on the right-hand side:
    \[\int x^{100}e^x\,dx = \tfrac{1}{101}x^{101}e^x - \tfrac{1}{101}\int x^{101}e^x\,dx \]
    
    If we let $u=x^{100}$, then we get a slightly simpler integral on the right-hand side:
    \[\int x^{100}e^x\,dx = x^{100}e^x - 100\int x^{99}e^x\,dx \]
    Every time we apply the integration by parts formula, the exponent on $x$ will go down by 1 in the integral on the right-hand side, until we eventually just have a multiple of $\int e^x \,dx$, which we know how to find. So we need to apply the integration by parts formula 100 times.
\end{solution}

\item 

For each of the following, say what the best choice for $u$ in the integration by parts formula would be.
\begin{enumerate}
    \item 
    \begin{problem}
    $\displaystyle \int x^{100} \ln(x)\,dx$
    \end{problem}
    \pspace{\stretch{0.3}}

    \begin{solution}
        The best choice is $u=\ln x$, which would lead to
        \[\int x^{100}\ln(x)\,dx = \tfrac{1}{101}x^{101}\ln(x) - \tfrac{1}{101}\!\int x^{100}\,dx \]
        and the integral on the right-hand side is simpler than the one we started with.

        If we used $u=x^{100}$ instead, we would get 
        \[\int x^{100}\ln(x)\,dx = x^{100}(x\ln x-x) - 100\int x^{99}(x\ln x - x) \,dx \]
        and the integral on the right-hand side is not simpler than what we started with.
    \end{solution}
    
    \item 
    \begin{problem}
    $\displaystyle \int x^{100} \cos(x)\,dx$
    \end{problem}
    \pspace{\stretch{0.3}}

        \begin{solution}
        The best choice is $u=x^{100}$, which would lead to
        \[\int x^{100}\cos(x)\,dx = x^{100}\sin(x) - 100\!\int x^{99}\sin (x)\,dx \]
        and the integral on the right-hand side is simpler than the one we started with.

        If we used $u=\cos(x)$ instead, we would get 
        \[\int x^{100}\cos(x)\,dx = -\tfrac{1}{101}x^{101}\sin(x) + \tfrac{1}{101}\int x^{101}\sin(x) \,dx \]
        and the integral on the right-hand side is not simpler than what we started with.
    \end{solution}
    
    \item 
    \begin{problem}
    $\displaystyle \int x^{100} \arctan(x)\,dx$
    \end{problem}
    \pspace{\stretch{0.3}}

            \begin{solution}
        The best choice is $u=\arctan x$, which would lead to
        \[\int x^{100}\arctan(x)\,dx = \tfrac{1}{101}x^{101}\arctan(x) - \tfrac{1}{101}\int \frac{x^{101}}{1+x^2} \,dx \]
        and the integral on the right-hand side is simpler than the one we started with.

        If we used $u=x^{100}$ instead, we would get 
        \[\int x^{100}\arctan(x)\,dx = x^{100}(x\arctan(x) - \tfrac{1}{2}\ln|1+x^2|) - 100\!\int x^{99}\left(x\arctan(x) - \tfrac{1}{2}\ln|1+x^2|\right)\,dx \]
        and the integral on the right-hand side is not simpler than what we started with.
    \end{solution}
\end{enumerate}


\begin{solution}
    
\end{solution}

\end{enumerate}

\psreflection{
Optional reading for next class is \S 29.1 in Gottlieb.
}


\end{document}